\documentclass[11pt]{article}
\usepackage{tabularx}
\usepackage[legalpaper, landscape, margin=1in]{geometry}
\usepackage[table]{xcolor} % Required for cell coloring
\usepackage[sfdefault]{atkinson}


\begin{document}
\pagenumbering{gobble}

\begin{tabularx}{\textwidth}{|c|X|X|X|}
    \hline
    \textbf{Week} & \textbf{Monday} & \textbf{Wednesday} & \textbf{Friday} \\
    \hline
    Week 1 & 
    
    01/13/25 \newline
    \textbf{Course overview and early works of sci-fi} \newline
    Critique of the first sci-fi film \newline
    Imagining human spaceflight in 1902\newline
    \textcolor{red}{\textbf{Science knowledge survey}} &
    
    01/15/25 \newline
    \newline
    Reference frames, acceleration and \newline
    artificial gravity &

    01/17/25 \newline
    \textbf{What is the nature of space and time?} \newline
    Classical physics and Newton's laws \newline
    \newline
    \textcolor{red}{\textbf{Quiz 1: Intro material (20 pts)}} \\


    \hline
    Week 2 & 

    \cellcolor{gray!20}01/20/25 \newline
    \newline
    \textbf{No Class: Martin Luther King Jr Day} \newline
    \newline
     &
    
    01/22/25 \newline
    \newline
    Newton's laws continued \newline
    Einstein and modern physics\ &

    01/24/25  \newline
    \newline
    Einstein and modern physics \newline
    Relativistic spacetime \newline
    \textcolor{red}{\textbf{Problem 1: acceleration}} \\


    \hline
    Week 3 & 

    01/27/25 \newline
    \newline
    Einstein and modern physics \newline
    Special relativity and time travel &
    
    01/29/25 \newline
    \newline
    General relativity \newline
    Black holes and gravitational waves&

    01/31/25  \newline
    \newline
    General relativity continued \newline
    Gravitational waves and the warp drive\newline
    \textcolor{red}{\textbf{Problem 2: time dilation}}\\
    
    \hline
    Week 4 & 
    
    02/03/25 \newline
    \newline
    "Future physics": Quantum gravity, worm   \newline
    holes, and the Multiverse Hypothesis \newline
    \textcolor{red}{\textbf{Exploration Paper 1: Space \& Time}} &
    
    \cellcolor{gray!20}02/05/25 \newline
    \newline
    \textbf{No Class: Spacetime Team Project} \newline
    \newline
    &

    \cellcolor{gray!20}02/05/25 \newline
    \newline
    \textbf{No Class: Spacetime Team Project}\newline\newline    
    \textcolor{red}{\textbf{Problem 3: General relativity}} \\


    \hline
    Week 5 & 
    
    02/10/25 \newline
    \textbf{What is the Universe made of?}\newline
    The standard model of particle physics \newline\newline
    \textcolor{red}{\textbf{Spacetime Team Project Due}} &
    
    02/12/25 \newline  
    \newline
    Dark matter and dark energy\newline
    \newline
    \textcolor{red}{\textbf{Quiz 2: Space and Time (85 pts)}} &

    02/14/25 \newline
    \newline
    Nuclei, atoms, and gases \\
    
    \hline
    Week 6 & 
    
    02/17/25 \newline
    \newline
    Liquids, solids, phase transitions \newline\newline
    &
    
    02/19/25 \newline  
    \newline
    Optical properties; transparent materials    
    &

    \cellcolor{gray!20}02/21/25 \newline
    \newline
    \textbf{No class today: Optical Team Project} \newline\newline
    \textcolor{red}{\textbf{Problem 4: matter \& materials}} \\

    \hline
    Week 7 & 
    
    02/24/25 \newline
    \newline
    Metamaterials and invisibility cloaks \newline\newline
    \textcolor{red}{\textbf{Optical Team Project due}}&
    
    02/26/25 \newline  
    \newline
    Chemical energy\newline\newline   
    \textcolor{red}{\textbf{Midterm Film Critique due}}&

    02/28/25 \newline
    \newline
    Nuclear weapons; matter \& antimatter \newline\newline
    \textcolor{red}{\textbf{Exploration Paper 2: Matter \& Energy}} \\
    \hline
\end{tabularx}
 
\begin{tabularx}{\textwidth}{|c|X|X|X|}

    \hline
    Week 8 & 
    
    \cellcolor{gray!20}03/03/25 \newline
    \textbf{No Class: Fall Break} 
    &
    
    \cellcolor{gray!20}03/05/25 \newline
    \textbf{No Class: Fall Break} 
    &

    \cellcolor{gray!20}03/07/25 \newline
    \textbf{No Class: Fall Break} 
    \\

    \hline
    Week 9 & 
    
    03/10/25 \newline
    \textbf{Can a machine become self aware?}\newline
    Information storage \& data processing \newline\newline
    &
    
    03/12/25 \newline
    \newline
    Intelligent robots \newline\newline
    \textcolor{red}{\textbf{Quiz 3: Matter and Energy (85 pts)}} &

    03/14/25 \newline
    \newline
    Bonus class day; Topic TBA \newline\newline
    \\
    \hline
    Week 10 & 
    
    03/17/25 \newline
    \newline
    The Turing Test \newline\newline
    &
    
    03/19/25 \newline
    \newline
    Asimov's Laws and robot behavior \newline\newline
    &

    03/21/25 \newline
    \newline
    Artificial conciousness \& robot ethics \newline
    
    \textcolor{red}{\textbf{Exploration Paper 3: AI}}\\

    \hline
    Week 11 & 
    
    03/24/25 \newline
    \textbf{Are we alone in the Universe?}\newline
    Conditions necessary for life to exist\newline
    \newline
    \textcolor{red}{\textbf{Quiz 4: AI (50 pts)}}&
    
    03/26/25 \newline
    \newline
    Finding habitable planets\newline
    \newline
    &

    03/28/25 \newline
    \newline
    \newline
    The science behind the SETI project\newline
    \textcolor{red}{\textbf{Exploration Paper 4: Extraterrestrial Intelligence}}\\

    \hline
    Week 12 & 
    
    03/31/25 \newline
    \textbf{What does it mean to be human?}\newline
    Insight from Frankenstein \newline
    Medical technology \newline
    \textcolor{red}{\textbf{Quiz 5: Extraterrestrial Intelligence (30 pts)}}&
    
    04/02/25 \newline
    \newline
    Cell structure and radiation damage \newline
    DNA, genome sequencing \& gene editing\newline
    &

    04/04/25 \newline
    \newline
    Human teleportation \newline
    \\
    
    \hline
    Week 13 & 

    04/07/25 \newline
    \newline
    What can we learn from an android \newline
    about what it means to be human?\newline
    \textcolor{red}{\textbf{Exploration Paper 5: Being Human}} &
    
    04/09/25 \newline
    \textbf{How can we solve our problems}\newline
    Public perception of science \newline
    The scientific method\newline &

    04/11/25 \newline
    \newline
    Misplaced faith in science \newline
    Misunderstanding of science\newline
    \textcolor{red}{\textbf{Quiz 6: Being Human (50 pts)}} \\

    \hline
    Week 14 & 
    
    04/14/25 \newline
    \newline
    Misuse of science in advertising \newline
    \textcolor{red}{\textbf{Final Film Critique due}} &
    
    \cellcolor{gray!20}04/16/25 \newline
    \newline
    \textbf{No Class: Thanksgiving Break} &

    \cellcolor{gray!20}04/18/25 \newline
    \newline
    \textbf{No Class: Thanksgiving Break} \\


    \hline
    Week 15 & 
    
    04/21/25 \newline
    \textbf{What lies ahead?} \newline
    Examples of sci-fi predicting the future \newline
    How sci-fi storytelling changes with time &
    
    04/23/25 \newline
    \newline
    Imagining the future \newline
    \textcolor{red}{\textbf{Quiz 7: Problem solving (30 pts)}}\newline &


    04/25/25 \newline
    \newline
    Last day of class -- closing discussion \\
    \hline
\end{tabularx}

\end{document}